\documentclass[11pt,a4paper]{article}
\usepackage[T1]{fontenc}
\usepackage{lmodern}
\usepackage[margin=23mm]{geometry}
\usepackage{amsmath,amssymb,booktabs,microtype}
\usepackage[hidelinks,hypertexnames=false]{hyperref}
\hypersetup{pdftitle={A first local hydrogen NLTE model},pdfauthor={exojax-escape}}
\pdfinfoomitdate=1
\pdftrailerid{}
\pdfsuppressptexinfo=15
\newcommand{\nH}{n_{\mathrm H}}
\newcommand{\kb}{k_{\mathrm B}}
\newcommand{\wn}{\widetilde\nu}
\title{A first local hydrogen NLTE model}
\author{exojax-escape}
\date{Definition revision 1 --- 20 November 2026}
\begin{document}
\maketitle

\section{What is solved}
The new model solves absolute hydrogen populations and the ionized fraction
at specified gas temperature $T$, nuclear density $\nH$ and radiation field
$J_\nu$. It resolves $1s,2s,2p,3s,3p,3d$, protons and electrons. The electron
density is determined by charge conservation, rather than fitted independently.
This is non-LTE: the level fractions are determined by transition rates, not
imposed Boltzmann or Saha populations. It is a local, prescribed-field
approximation, not a self-consistent atmosphere or a reproduction of the
Cloudy/sunbather calculation in Zhang et al.~\cite{zhang}.

The calculation follows this sequence:
\begin{align*}
 (T,\nH,J_\nu)&\ \longrightarrow\ \text{atomic rates}
 \ \longrightarrow\ \text{local rate balance and conservation},\\
 (n_{1s},\ldots,n_{3d},n_p)&\ \longrightarrow\
 \text{three H$\alpha$ branch weights}
 \ \longrightarrow\ \text{existing observer-ray transfer}.
\end{align*}
The final observer-directed transfer does not determine the radiation field
that pumps the atom. Shielding, Ly$\alpha$ trapping, internal recombination
radiation and radiation feedback are deferred.

The old isolated two-level closure and its records remain unchanged.
\path{doc/model.tex} and its PDF describe that historical baseline and the
original development proposal. This document defines the new model; its
machine-readable definitions are in
\path{mceg/contracts/escape.hydrogen_nlte.yaml}.

\section{The adopted atom and rates}
All densities are in cm$^{-3}$, frequencies $\nu$ in Hz, temperatures in K,
per-particle rates in s$^{-1}$ and two-body coefficients in cm$^3$\,s$^{-1}$.
Three-body coefficients are in cm$^6$\,s$^{-1}$. A rate matrix entry $P_{ij}$
means a transition \emph{from} $i$ \emph{to} $j$; its diagonal is zero.
The statistical weights in level order are $(2,2,6,2,6,10)$. Statistical
populations are assumed only within a fine-structure multiplet, not between
these six resolved states.

\paragraph{Bound--bound processes.}
Fine-structure energies and Einstein coefficients are aggregated from
Lehmann et al.~\cite{lehmann}, Appendix B, Tables 5--6. The five allowed
one-photon branches are $2p\to1s$, $3p\to1s$, $3s\to2p$, $3p\to2s$ and
$3d\to2p$. Spontaneous emission, absorption and stimulated emission are
included. Collision strengths from Table 7 are summed over the fine states
and interpolated linearly in $\log T$ and $\log\Upsilon$. Upward and downward
coefficients use the same adopted energy gaps and detailed balance.

Electron and proton $2s\to2p$ mixing use the four terms in Huang et
al.~\cite{huang}, Eqs.~12--13:
\begin{align*}
 q_{2s,2p}={}&1.394\times10^{-3}T^{-0.4055}e^{-13.35/T}
 +8.723\times10^{-4}T^{-0.3803}e^{-23.97/T}\\
 &+3.706\times10^{-3}T^{-0.2509}e^{-188.6/T}
 +3.654\times10^{-4}T^{-0.07835}e^{-1080/T}.
\end{align*}
The first two terms are electron collisions and the last two proton
collisions. They can be combined because this pure-H model has $n_e=n_p$.
The reverse coefficient follows detailed balance with the adopted small
$2s/2p$ energy difference. Mixing between $n=3$ sublevels is omitted.

The spontaneous $2s\to1s$ two-photon rate is $8.2$ s$^{-1}$. Induced and
inverse two-photon processes are absent. Consequently, the \emph{complete}
atom is not asserted to reach exact Planck--Saha LTE at arbitrary density.
Detailed-balance checks for the included one-photon and collisional processes
remain useful and are narrower statements.

\paragraph{Bound--free processes.}
Photoionization includes all six levels, including the Balmer continuum.
The Kramers approximation uses bound-free Gaunt factor one and the same
cross section for each orbital $l$ at fixed principal quantum number.
Milne recombination uses the same cross section and includes stimulated
recombination. The default quadrature has 128 Gauss--Legendre nodes per
level, uniform in the quadrature variable $\log\nu$, from its threshold
to a fixed photon energy of 100 eV. The public rate API also accepts
prescribed continuum nodes, weights and occupations. Harder spectra require
an appropriate integration range and a new convergence check.

Electron-ionization coefficients use $q^{\rm ion}=aT^b\exp(c/T)$ from
Garc\'ia Mu\~noz and Schneider~\cite{garcia}, Appendix, Table 2:
\begin{center}
\begin{tabular}{@{}lrrr@{}}
\toprule
Principal level & $a$ & $b$ & $c$ [K]\\
\midrule
1 & $1.51$ & $-2.09$ & $-1.68\times10^5$\\
2 & $1.35\times10^{-7}$ & $-0.033$ & $-3.89\times10^4$\\
3 & $1.96\times10^{-7}$ & $0.15$ & $-1.85\times10^4$\\
\bottomrule
\end{tabular}
\end{center}
With $T$ in K, $a$ carries the units needed to give cm$^3$\,s$^{-1}$.
Each principal-level fit is shared by its $l$ sublevels, including the
shell-averaged $n=3$ fit. The reported fit accuracy is about 20\% on
4000--14000 K; this does not characterize every other rate. Inverse
three-body recombination is set by the same truncated-atom Saha ratio.
No recombination into $n>3$, cascade correction or Case-B closure is added.

The supported \emph{input} temperature interval is the intersection of the
adopted data ranges, $0.5\,{\rm eV}/\kb\simeq5802.26$ K through 14000 K.
Inputs outside it produce invalid rates rather than extrapolated fits.
The interval is not a validated atmospheric applicability domain.

\section{Solving and differentiating the populations}
At a trial electron density, the six bound-level equations are linear in
$b_i=n_i/n_p$. A $6\times6$ solve gives $b$; the only remaining unknown is
$z=\log(n_e/\nH)$. Nuclear conservation supplies the scalar residual in
Eq.~\eqref{eq:escape.hydrogen_5f_nlte:charge_5f_root:1}.
The code uses a fixed bracket $-80\le z\le0$ and 80 bisection iterations by
default. It requires at least one positive photoionization rate. Exactly
unirradiated or disconnected systems can have neutral absorbing branches
and lie outside this initial solver domain.

The solver returns the residual, fractional conservation error, scaled
rate-balance error, root slope, input-validity flag and convergence flag.
The three error measures must be below $10^{-10}$ and the slope finite and
positive. Failed solutions return NaN populations; the flags are part of
the result, not optional scientific evidence. These are numerical convergence
criteria, not bounds on atomic or atmospheric error.

JAX's \path{lax.custom_root} differentiates the same residual, including the
matrix solve and the feedback of $n_e$ into collision and recombination
rates. It does not differentiate the bisection branch choices. A finite,
nonzero local Jacobian and a sufficiently converged root are prerequisites.
Gradient comparisons must be supported by the saved numerical receipt;
recording the formula alone does not validate derivatives or a NUTS analysis.

Local equilibrium omits population transport. Comparing a relaxation time
with $r/v$ can provide a preliminary diagnostic, but is not the full
population-gradient or radiation-feedback timescale test. Such diagnostics
do not establish the local-equilibrium approximation in every atmospheric cell.

\section{Connection to the spectrum and wind}
The new forward model uses a prescribed gas profile that starts at 6000 K
and approaches 10000 K by default. The continuum lower boundary remains
4500 K; gas and boundary brightness temperatures play different roles.
The initial radiation closure is $J_\nu=W B_\nu(T_{\rm rad})$, with default
$T_{\rm rad}=10000$ K and $W=0.025$. Those choices are inputs, not a solved
stellar or planetary spectrum.

The Parker velocity and spherical continuity determine $\nH$. The
nonisothermal wind uses the same prescribed gas profile but retains an
external constant mean molecular weight (default $\mu=0.5$). The population
solution supplies the diagnostic $\mu_{\rm NLTE}=1/(1+n_p/\nH)$; it does
not feed back into the wind or solve an energy equation. Optional tides
remain the existing one-dimensional linear effective-gravity approximation.

H$\alpha$ extinction and emissivity are summed from $3s\to2p$, $3p\to2s$
and $3d\to2p$, with each branch's actual lower and upper populations.
Replacing those populations by statistical shell fractions would lose the
resolved non-LTE information. The approximation is instead a common
rest center and normalized line profile; only the NIST center is supported.
The three branch centers implied by the adopted energies are approximately
15233.06043, 15233.33085 and 15233.17984 cm$^{-1}$, spanning 5.32 km\,s$^{-1}$.
One shared center omits their population-dependent centroid shift. The inherited
unresolved-line natural width also omits the branch-specific upper-plus-lower
lifetimes, whose summed decay rates are roughly 4--16 times larger than that
width prescription. Precise centroids and damping wings are not claimed.
For the existing positive-opacity formal solver, zero or negative integrated
extinction yields a NaN spectrum. No opacity clipping substitutes for a
maser-capable transfer calculation.

The inherited disk geometry, observer-ray transfer, continuum boundary,
broadening and velocity prescriptions retain their earlier limitations.
In particular, this integration adds neither an internal pumping-field solve
nor an observational reduction model.

\section{Evidence and model choices}
The MCEG contract separates numerical checks, formula/code correspondence
and physical adequacy. Decisions record the six-state truncation, external
radiation, atomic approximations, temperature restriction and branch sum,
together with limitations and reconsideration conditions. An unknown physical
Assessment can coexist with an explicitly adopted approximation.

The equation appendix below is generated from that contract and has named
Bindings to Python source. MCEG can check the recorded definition identities,
source declarations and saved test artifacts; it does not prove physical
correctness. The actual receipt determines which numerical claims are
supported. The retained example includes a separate synthetic-spectrum NUTS
smoke run: one chain, 32 warmup steps and 32 retained samples. Finite samples
and no reported divergence confirm that this small software path executed;
they do not establish posterior convergence or parameter recovery. No Cloudy
comparison or observed-spectrum fit is implied.
The old two-level contract, source snapshot and six-test receipt are preserved.
Neither MCEG nor a TeX engine is an \texttt{escape} runtime dependency.

\appendix
\section{Recorded equations}\label{sec:equations}
The following definitions are generated from the schema-0.4 contract. Their
symbols and scope remain inspectable alongside the code bindings.
% Generated by mceg equations; edit the contract, then regenerate.
% Requires amsmath. Authored math is retained; TeX is not executed here.
% Definitions and hashes do not establish scientific validity.

% equation-sha256: c42f5f5133532de6abdf4b0ca494f579910cc56e1623ee9efcd61a79fbae387f
\paragraph{Equation level\_rates (revision 1)}
Allowed bound-bound rates, with prescribed photon occupation and Maxwellian collisions.
\begin{equation}
\begin{aligned}
P_{lu}=n_e q_{lu}+\frac{g_u}{g_l}A_{ul}\bar n_{ul},\qquad P_{ul}=n_e q_{ul}+A_{ul}(1+\bar n_{ul}),\\\bar n_{ul}=\frac{c^2\overline J_{ul}}{2h\nu_{ul}^3}
\end{aligned}
\label{eq:escape.hydrogen_5f_nlte:level_5f_rates:1}
\end{equation}
\noindent\textbf{Scope:} One-photon branches only; the one-way 2s two-photon decay is added separately.

\begin{description}
\item[$P_{ij}$] Total rate from bound level i to j (s\textasciicircum{}-1; ID: p).
\item[$\bar n_{ul}$] Profile-averaged photon occupation (1; ID: occupation).
\item[$q_{ij}$] Collision coefficient; proton mixing uses ne=np (cm\textasciicircum{}3 s\textasciicircum{}-1; ID: q).
\end{description}

\noindent\textbf{Assumptions:}
\begin{itemize}
\item No radiation feedback from the solved populations.
\item Bound-bound arrays are indexed [from, to], with zero diagonal.
\end{itemize}

% equation-sha256: 8e82e41de0018aeecd9cc9fe48380c10ece74fd0c9a82706434f99d4e7fcd700
\paragraph{Equation collision\_balance (revision 1)}
Electron excitation and de-excitation share an interpolated effective collision strength.
\begin{equation}
\begin{aligned}
q_{lu}=\frac{8.629\times10^{-6}\Upsilon_{lu}(T)}{g_l\sqrt{T}}\exp\!\left(-\frac{\Delta E_{ul}}{k_{\rm B}T}\right),\qquad q_{ul}=\frac{8.629\times10^{-6}\Upsilon_{lu}(T)}{g_u\sqrt{T}}
\end{aligned}
\label{eq:escape.hydrogen_5f_nlte:collision_5f_balance:1}
\end{equation}
\noindent\textbf{Scope:} The adopted temperature table; mixing between degenerate 2s and 2p has its own prescription.

\begin{description}
\item[$T$] Electron temperature (K; ID: temperature).
\item[$\Upsilon_{lu}$] Dimensionless thermally averaged collision strength (1; ID: upsilon).
\item[$q_{ij}$] Maxwellian electron collision coefficient (cm\textasciicircum{}3 s\textasciicircum{}-1; ID: q).
\item[$\mathcal R$] Three returned collision-coefficient arrays (mixed cgs; ID: rates).
\end{description}

\noindent\textbf{Assumptions:}
\begin{itemize}
\item Boltzmann detailed balance constrains the inverse rate, not the accuracy of the adopted table.
\end{itemize}

% equation-sha256: de870e3d3c2c5d1f4819a05bbcb7359de89c0686810dc2f0da227a79746381ee
\paragraph{Equation bound\_free (revision 1)}
Photoionization, Milne radiative recombination and inverse three-body recombination.
\begin{equation}
\begin{aligned}
K_{ik}=\frac{8\pi\nu_{ik}^2}{c^2}\sigma_i(\nu_{ik}),\quad \Gamma_i=\sum_k w_{ik}K_{ik}\bar n_{ik},\\\alpha_i=\frac{1}{S_i(T)}\sum_k w_{ik}K_{ik}(1+\bar n_{ik})\exp\!\left(-\frac{h\nu_{ik}}{k_{\rm B}T}\right),\\S_i(T)=\frac{2}{g_i}\left(\frac{2\pi m_e k_{\rm B}T}{h^2}\right)^{3/2}\exp\!\left(-\frac{\chi_i}{k_{\rm B}T}\right),\quad q_i^{(3)}=q_i^{\rm ion}/S_i(T)
\end{aligned}
\label{eq:escape.hydrogen_5f_nlte:bound_5f_free:1}
\end{equation}
\noindent\textbf{Scope:} The same approximate cross section is used in both radiative directions, including stimulated recombination.

\begin{description}
\item[$\Gamma_i$] Photoionization rate from level i (s\textasciicircum{}-1; ID: gamma).
\item[$\alpha_i$] Radiative recombination coefficient into level i (cm\textasciicircum{}3 s\textasciicircum{}-1; ID: alpha).
\item[$S_i$] Truncated-atom Saha ratio ne*np/ni (cm\textasciicircum{}-3; ID: saha).
\item[$q_i^{(3)}$] Inverse three-body coefficient (cm\textasciicircum{}6 s\textasciicircum{}-1; ID: q3).
\item[$T$] Electron temperature (K; ID: temperature).
\item[$\bar n_{ik}$] Prescribed continuum photon occupation (1; ID: occupation).
\item[$\nu_{ik}$] Continuum quadrature frequencies for level i (Hz; ID: frequency).
\item[$w_{ik}$] Quadrature integration weights for dnu (Hz; ID: weights).
\end{description}

\noindent\textbf{Assumptions:}
\begin{itemize}
\item Hydrogenic Kramers cross sections with bound-free Gaunt factor one.
\item No recombination or cascades through omitted n>3 levels.
\item Default quadrature: 128 Gauss-Legendre nodes in log frequency per threshold through 100 eV. Harder spectra need a separate convergence check.
\end{itemize}

% equation-sha256: fe1aca8ccedb9f80d3f7ed08cbbcafcebfad5932e8cd6755545fba88ac2dba21
\paragraph{Equation atomic\_closure (revision 1)}
Kramers photoionization and the adopted electron-ionization fit.
\begin{equation}
\begin{aligned}
\sigma_i(\nu)=\frac{2.815\times10^{29}}{N_i^{5}\nu^3}\,\mathbf 1_{h\nu\ge\chi_i},\qquad q_i^{\rm ion}=a_i T^{b_i}\exp(c_i/T)
\end{aligned}
\label{eq:escape.hydrogen_5f_nlte:atomic_5f_closure:1}
\end{equation}
\noindent\textbf{Scope:} N\_i denotes the principal quantum number of bound level i; T is restricted to the adopted rate domain.

\begin{description}
\item[$\nu$] Photon frequency (Hz; ID: frequency).
\item[$\sigma_i$] Gaunt-one bound-free cross section (cm\textasciicircum{}2; ID: sigma).
\item[$T$] Electron temperature (K; ID: temperature).
\item[$N_i$] Principal quantum number of bound level i (1; ID: principal).
\end{description}

\noindent\textbf{Assumptions:}
\begin{itemize}
\item Cross sections are independent of orbital l at fixed n.
\item The fit coefficients and their units are tabulated in the companion document; the adopted data domain is restricted.
\item Ionization fits are shared by all orbital-l levels of each principal shell, including the shell-averaged n=3 fit.
\end{itemize}

% equation-sha256: 97820a31cf2d375b0f1c532c806d960cc3b1d14c19c0a97efce3ec80309651f7
\paragraph{Equation level\_balance (revision 1)}
Local statistical equilibrium, hydrogen-nucleus conservation and pure-H charge conservation.
\begin{equation}
\begin{aligned}
Q_i=\sum_{j\ne i}(n_jP_{ji}-n_iP_{ij})+n_e^2\alpha_i+n_e^3q_i^{(3)}-n_i(\Gamma_i+n_eq_i^{\rm ion})=0,\\\sum_i n_i+n_p=n_{\rm H},\qquad n_e=n_p
\end{aligned}
\label{eq:escape.hydrogen_5f_nlte:level_5f_balance:1}
\end{equation}
\noindent\textbf{Scope:} Six bound nl states plus one proton/electron density; no transport or energy equation.

\begin{description}
\item[$n_{\rm H}$] Prescribed total hydrogen-nucleus density (cm\textasciicircum{}-3; ID: density).
\item[$n_i$] Bound-level number density (cm\textasciicircum{}-3; ID: bound).
\item[$n_e=n_p$] Electron and proton number density (cm\textasciicircum{}-3; ID: electron).
\item[$Q_i$] Net production of bound level i (cm\textasciicircum{}-3 s\textasciicircum{}-1; ID: net).
\end{description}

\noindent\textbf{Assumptions:}
\begin{itemize}
\item Local reaction equilibrium neglects advection; its applicability to each outflow cell remains unvalidated.
\end{itemize}

% equation-sha256: 34ed2e8256a19fc714f21ee0fbfc1d762ff525ecb0bf5ffeeb395908d75d6655
\paragraph{Equation fixed\_electron\_system (revision 1)}
At fixed electron density the bound-level problem is one linear solve.
\begin{equation}
\begin{aligned}
M_{ij}=\delta_{ij}\!\left(\sum_kP_{ik}+\Gamma_i+n_eq_i^{\rm ion}\right)-P_{ji},\\M(n_e)b=n_e\alpha+n_e^2q^{(3)},\qquad b_i=n_i/n_e
\end{aligned}
\label{eq:escape.hydrogen_5f_nlte:fixed_5f_electron_5f_system:1}
\end{equation}
\noindent\textbf{Scope:} Nonnegative rates with zero bound-bound diagonals and a nonsingular loss matrix.

\begin{description}
\item[$n_e$] Trial electron density (cm\textasciicircum{}-3; ID: electron).
\item[$\mathcal R$] RateCoefficients tuple; units are recorded for each field (mixed cgs; ID: rates).
\item[$b$] Bound-level density divided by proton density (1; ID: ratios).
\item[$M$] Bound-level loss matrix (s\textasciicircum{}-1; ID: matrix).
\end{description}

\noindent\textbf{Assumptions:}
\begin{itemize}
\item Rows are divided by their total outgoing rate for numerical scaling; the equations are unchanged.
\end{itemize}

% equation-sha256: 3f5180c937a92ff38419de35172f123992429a0b21ca60e2a687fcb3a5e00c5b
\paragraph{Equation charge\_root (revision 1)}
The remaining scalar root enforces nuclear conservation without a redundant proton equation.
\begin{equation}
\begin{aligned}
z=\log(n_e/n_{\rm H}),\qquad F(z)=z+\log\!\left(1+\sum_i b_i(n_{\rm H}e^z)\right)=0,\\n_e=n_p=n_{\rm H}e^z,\qquad n_i=n_e b_i
\end{aligned}
\label{eq:escape.hydrogen_5f_nlte:charge_5f_root:1}
\end{equation}
\noindent\textbf{Scope:} One cell, float64, positive nH; at least one positive photoionization rate; bracket -80 <= z <= 0.

\begin{description}
\item[$n_{\rm H}$] Prescribed total hydrogen-nucleus density (cm\textasciicircum{}-3; ID: density).
\item[$\mathcal R$] RateCoefficients tuple supplied to the solver (mixed cgs; ID: rates).
\item[$z$] Log ionized fraction (1; ID: z).
\item[$F$] Scalar conservation residual (1; ID: residual).
\end{description}

\noindent\textbf{Assumptions:}
\begin{itemize}
\item Eighty bisection iterations by default.
\item Invalid or unconverged solutions return NaN populations and explicit diagnostics.
\end{itemize}

% equation-sha256: 22aaa52f5fa1cde7fbc7a20097500a389458be35c7d3bac71dfd056435a31ab1
\paragraph{Equation implicit\_derivative (revision 1)}
Implicit differentiation uses the same residual as the converged root.
\begin{equation}
\begin{aligned}
F_z\,\delta z=-F_\theta\,\delta\theta,\qquad \frac{dz}{d\theta}=-\frac{F_\theta}{F_z}
\end{aligned}
\label{eq:escape.hydrogen_5f_nlte:implicit_5f_derivative:1}
\end{equation}
\noindent\textbf{Scope:} A locally isolated converged root; derivatives include electron-density feedback and the linear solve.

\begin{description}
\item[$z$] Log ionized fraction (1; ID: z).
\item[$\theta$] Density, rate coefficients and their upstream parameters (parameter dependent; ID: parameters).
\item[$F$] Scalar conservation residual (1; ID: residual).
\end{description}

\noindent\textbf{Assumptions:}
\begin{itemize}
\item jax.lax.custom\_root differentiates the equations, not the bisection decisions.
\item A positive finite root slope is required by the implemented convergence diagnostics.
\end{itemize}

% equation-sha256: f04440a43051e18fe765f2c028e85bde2f27847cff31be78518b481101c53499
\paragraph{Equation branch\_coefficients (revision 1)}
Sum the three H-alpha branches before multiplying by a shared normalized profile.
\begin{equation}
\begin{aligned}
\mathcal A=\sum_{u\to l}\frac{A_{ul}}{8\pi c\widetilde\nu_0^2}\left(\frac{g_u}{g_l}n_l-n_u\right),\qquad \mathcal E=\sum_{u\to l}\frac{hc\widetilde\nu_0 A_{ul}}{4\pi}n_u,\\\chi_{\widetilde\nu}=\mathcal A\phi_{\widetilde\nu},\qquad j_{\widetilde\nu}=\mathcal E\phi_{\widetilde\nu},\qquad \int\phi_{\widetilde\nu}\,d\widetilde\nu=1
\end{aligned}
\label{eq:escape.hydrogen_5f_nlte:branch_5f_coefficients:1}
\end{equation}
\noindent\textbf{Scope:} Branches 3s->2p, 3p->2s and 3d->2p use their actual nl populations; only the common NIST center is supported.

\begin{description}
\item[$n_i$] Actual nl-level number densities (cm\textasciicircum{}-3; ID: bound).
\item[$\widetilde\nu_0$] Common H-alpha rest wavenumber (cm\textasciicircum{}-1; ID: center).
\item[$\mathcal A$] Extinction integrated over wavenumber (cm\textasciicircum{}-2; ID: alpha).
\item[$\mathcal E$] Emissivity integrated over wavenumber (erg s\textasciicircum{}-1 cm\textasciicircum{}-3 sr\textasciicircum{}-1; ID: emission).
\end{description}

\noindent\textbf{Assumptions:}
\begin{itemize}
\item No statistical redistribution between nl levels is made.
\item Common-profile line transfer is evaluated only for positive integrated extinction; zero or negative extinction is reported through NaN spectra.
\item The shared natural width omits branch-specific upper-plus-lower lifetimes; no precise centroid or wing accuracy is claimed.
\end{itemize}

% equation-sha256: 752f9d705a4ef1597418e6662ab6bc8f33e389b0feaef9e6afb5a124fa7f8d4d
\paragraph{Equation relaxation\_time (revision 1)}
Slow local kinetic relaxation includes the charge response to changing bound populations.
\begin{equation}
\begin{aligned}
K_{ij}=\left.\frac{\partial Q_i}{\partial n_j}\right|_{n_{\rm H},T,J_\nu},\quad n_e=n_p=n_{\rm H}-\sum_i n_i,\\\tau_{\rm relax}=-\frac{1}{\max_a\operatorname{Re}\lambda_a(K)}\quad\text{if all }\operatorname{Re}\lambda_a(K)<0
\end{aligned}
\label{eq:escape.hydrogen_5f_nlte:relaxation_5f_time:1}
\end{equation}
\noindent\textbf{Scope:} Optional diagnostic evaluated at converged local populations; no radiation or transport feedback.

\begin{description}
\item[$n_{\rm H}$] Fixed total hydrogen-nucleus density (cm\textasciicircum{}-3; ID: density).
\item[$n_i$] Independent bound-level number densities (cm\textasciicircum{}-3; ID: bound).
\item[$\mathcal R$] The prescribed RateCoefficients used by the population solve (mixed cgs; ID: rates).
\item[$\tau_{\rm relax}$] Slowest local damping time (s; ID: time).
\end{description}

\noindent\textbf{Assumptions:}
\begin{itemize}
\item The reduced kinetic Jacobian has no redundant nuclear-conservation mode.
\item Nondecaying or unstable modes return infinity; invalid or nonfinite inputs return NaN.
\item Comparison with r/v is a characteristic diagnostic, not a proof that the relevant population-gradient transport time is long.
\end{itemize}


\begin{thebibliography}{9}
\bibitem{lehmann}
Lehmann et al. (2020), Appendix B, Tables 5--7 and Eq.~49.
\url{https://arxiv.org/abs/2010.01042v1}.
\bibitem{garcia}
A. Garc\'ia Mu\~noz and P. C. Schneider (2019),
\emph{Rapid escape of ultra-hot exoplanet atmospheres driven by hydrogen
Balmer absorption}, Appendix and Table 2.
\url{https://arxiv.org/abs/1910.00267v1}.
\bibitem{huang}
Huang et al. (2023), Eqs.~12--13.
\url{https://arxiv.org/abs/2304.07352v1}.
\bibitem{zhang}
Y. Zhang et al. (2026), \emph{Hydrogen airglow from an escaping ultrahot
Jupiter atmosphere}, Sections IV--V.
\url{https://arxiv.org/abs/2606.23899v1}.
\end{thebibliography}
\end{document}
